% -*- mode:LaTex; mode:visual-line; mode:flyspell; fill-column:75-*-

\chapter{Data Collection and 2D Disease Detection} \label{chapDataset}

Realizing the active perception approach described in Chapter~\ref{chapIntroduction} first requires a perception system that can recognize fire blight symptoms in close-range, dormant-season imagery.
This chapter develops that system from the ground up: the handheld multi-modal camera rig used to collect orchard imagery, the dataset assembled from multiple collection sessions across several rig revisions, the SAM-assisted and cross-modal annotation pipeline used to label that data efficiently, and the 2D instance segmentation models trained on the resulting dataset.
The detector trained here is deployed as the perception front-end for the semantic mapping and next-best-view planning system presented in Chapter~\ref{chapNBV}.

\section{Sensing Platform} \label{secSensingPlatform}

Data collection was performed with a handheld, multi-modal camera rig rather than the autonomous manipulator platform used for inspection (Erwin), since rig development and dataset collection began before the manipulator-based system was operational and continued in parallel with it.
The rig combines a flash-illuminated stereo RGB camera, a 25-band NIR multispectral camera, and, in early versions, a wide-field RGB-D camera, all rigidly mounted to a common frame and faced in the same direction so that a single operator could walk the orchard rows and capture closely registered, multi-modal imagery of trunks and branches at close range.
The rig went through several physical and electrical revisions over the course of data collection, described in Section~\ref{secDesignIterations}; unless otherwise noted, the descriptions below refer to the final version used for the majority of dataset collection.

% TODO: Include a photo of the camera rig (full rig, with a blown-up
% inset/callout showing the individual camera mounts and connectors).
% \begin{figure}[h]
%   \centering
%   \includegraphics[width=0.85\textwidth]{figs/dataset/sensor_rig}
%   \caption{TODO: Handheld camera rig showing camera placement and mounting,
%            with an inset detail view of the camera mounts.}
%   \label{figSensorRig}
% \end{figure}

\subsection{Flash Stereo RGB Camera (FLIR Firefly)} \label{secFlirCamera}

The primary RGB sensor is a flash-illuminated stereo camera built around FLIR Firefly cameras, developed by Silwal et al.~\cite{Silwal2021FlashCamera} to provide consistent, shadow-free illumination of woody tissue regardless of ambient lighting conditions.
This is especially important for dormant-season fire blight inspection: cankers and shepherd's crooks are distinguished from healthy bark largely by subtle texture and color cues (sunken or darkened lesions, cracked bark) that are easily washed out or obscured by the variable, often low and oblique, natural lighting found in orchard rows.
The camera receives a hardware (serial) trigger command that synchronizes the flash pulse with the stereo image exposure, so that both the left and right frames are captured under identical, controlled illumination.
% TODO: Resolution, frame rate, lens focal length/FOV, stereo baseline,
% and effective depth range of the flash stereo pair.

\subsection{25-Band NIR Multispectral Camera (Ximea)} \label{secXimeaCamera}

A Ximea multispectral camera captures 25 narrowband channels in the near-infrared range, providing spectral information that is not visible in RGB but that has been shown to be discriminative for fire blight infection due to changes in tissue moisture content and pigmentation as cankers develop.
% TODO: Spectral range and per-band wavelengths; spatial resolution;
% snapshot mosaic vs. pushbroom acquisition; expected spectral signatures
% distinguishing healthy from infected tissue.
In early rig versions the Ximea camera was triggered independently by software, while later versions placed it on the same hardware trigger line as the flash stereo camera (Section~\ref{secDesignIterations}) so that NIR and RGB frames could be captured with tight, deterministic timing rather than being subject to software scheduling jitter.

\subsection{ZED2i RGB-D Camera} \label{secZED2i}

Early rig versions also included a Stereolabs ZED2i stereo RGB-D camera, intended to provide a wider-field RGB overview of the inspection area alongside depth for context and mapping.
% TODO: Frame rate, depth range, and FOV of the ZED2i as configured.
The ZED2i was ultimately removed from later rig versions: its imaging resolution was not fine enough to resolve the small-scale features needed to detect dormant fire blight symptoms, particularly bark texture on cankers and thin, small shepherd's crooks, so it contributed little to the dataset beyond redundant coarse RGB context already covered by the flash stereo camera.

\subsection{Platform and Mounting} \label{secPlatformMounting}

The three cameras are stacked vertically on the rig frame, all facing the same direction: the flash stereo camera at the top, the Ximea multispectral camera centered beneath it, and, in early versions, the ZED2i at the bottom.
A handle is mounted to the back of the rig for handheld operation in the field, alongside a touchscreen used to interface with the onboard computer.
In the field, the rig was driven by an NVIDIA Jetson Orin, powered by a drone battery pack for untethered operation away from line power.
% TODO: confirm exact Orin variant (AGX/NX/Nano) and module spec.
The Jetson connects to all cameras over USB and records synchronized data as ROS bags; incoming frames from each modality are time-synced on capture so that downstream annotation and label-transfer steps (Section~\ref{secCrossModal}) operate on temporally aligned, registered imagery.

\subsection{Design Iterations and Triggering} \label{secDesignIterations}

The rig went through several design revisions driven primarily by the practical demands of repeated field deployment and by the need for tighter timing synchronization across modalities.

\textbf{Camera loading and mounting.}
In the earliest version of the rig, all cameras were loaded into the base mounting block from the back, with the touchscreen then attached behind them for viewing images and triggering capture.
Because the rig was frequently disassembled and reassembled in the field, swapping cameras in and out of a back-loaded mount proved cumbersome, so later versions switched to an entirely front-loaded design, with slides added to lock each camera in place once positioned.
The spacing between cameras was also reduced in later versions to increase the overlap between their fields of view, improving the consistency of cross-modal label transfer between the flash RGB and NIR frames.

\textbf{Triggering and synchronization.}
The flash stereo camera was always driven by a hardware trigger, since the serial trigger command was needed to synchronize the flash pulse with stereo exposure.
The Ximea and ZED2i cameras, in contrast, were originally triggered by software, which introduced timing slop between frames of different modalities.
% TODO: quantify typical software-trigger timing offset, if known.
In later rig versions, the Ximea camera was moved onto the same hardware trigger line as the flash stereo camera, eliminating this slop between the two modalities retained in the final design; the ZED2i was dropped from the platform around the same time for the resolution limitations described in Section~\ref{secZED2i}.

\textbf{Control interface.}
The control software was ported from Python to C++ to improve capture and preview throughput, and the user interface was redesigned around large, glove-friendly buttons and sliders suited to outdoor field use.
The interface exposed direct control over flash stereo exposure, flash emitter time/intensity, and Ximea exposure time and gain, and the live preview was reduced from five camera views down to two (left flash stereo and Ximea NIR), the two modalities retained for the final dataset.

\section{Data Collection Protocol} \label{secDataCollection}

% TODO: Describe how data was collected.
% - Orchard location(s), dates, tree varieties.
% - Disease severity levels present.
% - Collection strategy: row traversal, distances to canopy, viewpoint angles.
% - Number of collection sessions / total coverage.
% - Environmental conditions (lighting, weather).

% \begin{figure}[h]
%   \centering
%   \includegraphics[width=\textwidth]{figs/dataset/orchard_map}
%   \caption{TODO: Map or photo of the orchard collection site
%            with collection paths overlaid.}
%   \label{figOrchardMap}
% \end{figure}

\section{Dataset Description} \label{secDatasetDescription}

% TODO: Summarize the dataset statistics.
% - Total number of image frames per modality.
% - Number of annotated images.
% - Class distribution (healthy vs. infected; parts: shoot, leaf, fruit, etc.).
% - Train / validation / test split.
% - Comparison to existing public datasets (if any).

% \begin{table}[h]
%   \centering
%   \caption{TODO: Dataset statistics table.}
%   \label{tabDatasetStats}
%   \begin{tabular}{lcccc}
%     \hline
%     Modality & Total frames & Annotated & Infected & Healthy \\
%     \hline
%     Flash stereo RGB & & & & \\
%     25-band NIR & & & & \\
%     ZED2i RGB & & & & \\
%     \hline
%   \end{tabular}
% \end{table}

\section{Annotation Pipeline} \label{secAnnotation}

% TODO: Describe how ground-truth labels were generated.

\subsection{Expert Manual Annotation} \label{secManualAnnotation}

% TODO: Describe how images were reviewed and labeled by domain experts
% (or the authors + expert verification). Tools used (LabelImg, CVAT, etc.).
% Inter-annotator agreement if applicable.

\subsection{SAM-Assisted Annotation} \label{secSAMAnnotation}

% TODO: Describe use of the Segment Anything Model (SAM) to accelerate
% annotation. Workflow: expert provides seed points → SAM generates masks →
% expert reviews/corrects. Quantify time savings.

% \begin{figure}[h]
%   \centering
%   \includegraphics[width=\textwidth]{figs/dataset/annotation_pipeline}
%   \caption{TODO: Annotation pipeline diagram showing SAM-assisted workflow.}
%   \label{figAnnotationPipeline}
% \end{figure}

\subsection{Cross-Modal Label Transfer} \label{secCrossModal}

% TODO: Describe how RGB annotations are transferred to NIR frames.
% - Camera calibration / extrinsics between FLIR and Ximea.
% - Projection-based label mapping.
% - Handling of occlusion, misalignment, and spectral differences.
% Include a figure showing paired RGB + NIR frames with transferred labels.

% \begin{figure}[h]
%   \centering
%   \includegraphics[width=\textwidth]{figs/dataset/cross_modal_transfer}
%   \caption{TODO: Example of cross-modal label transfer from RGB to NIR.}
%   \label{figCrossModal}
% \end{figure}

\section{2D Disease Detection Models} \label{sec2DDetection}

% TODO: Describe the detection models trained on this dataset.

\subsection{Model Architecture} \label{secModelArch}

% TODO: Describe the YOLO variant(s) used (e.g., YOLOv8, YOLOv11).
% - Input resolution, number of classes.
% - Pre-training (COCO weights) and fine-tuning procedure.
% - Separate models for RGB vs. NIR; or fused model.
% - Loss function, optimizer, augmentation strategy.

\subsection{Training Details} \label{secTraining}

% TODO: Hyperparameters, hardware, training duration.
% \begin{table}[h]
%   \centering
%   \caption{TODO: Training configuration.}
%   \label{tabTrainingConfig}
%   \begin{tabular}{ll}
%     \hline
%     Hyperparameter & Value \\
%     \hline
%     Architecture & \\
%     Input size & \\
%     Optimizer & \\
%     Learning rate & \\
%     Epochs & \\
%     Batch size & \\
%     \hline
%   \end{tabular}
% \end{table}

\subsection{Detection Results} \label{secDetectionResults}

% TODO: Report quantitative results (mAP@0.5, precision, recall, F1)
% on the test split. Compare RGB vs. NIR models; show per-class breakdown.
% Include precision-recall curves.

% \begin{table}[h]
%   \centering
%   \caption{TODO: Detection performance comparison.}
%   \label{tabDetectionResults}
%   \begin{tabular}{lcccc}
%     \hline
%     Model & mAP@0.5 & Precision & Recall & F1 \\
%     \hline
%     RGB (YOLOv?) & & & & \\
%     NIR (YOLOv?) & & & & \\
%     RGB+NIR fused & & & & \\
%     \hline
%   \end{tabular}
% \end{table}

% \begin{figure}[h]
%   \centering
%   \includegraphics[width=\textwidth]{figs/dataset/detection_examples}
%   \caption{TODO: Qualitative detection examples on test images
%            (RGB and NIR side by side with predicted bounding boxes).}
%   \label{figDetectionExamples}
% \end{figure}

\subsection{Failure Mode Analysis} \label{secFailureModes}

% TODO: Analyze where and why the detector fails.
% - Small or partially occluded lesions.
% - Lighting artifacts.
% - Early-stage infection with subtle appearance.
% Motivate downstream uncertainty-aware mapping and NBV planning.

\section{Summary} \label{secDatasetSummary}

% TODO: Summarize contributions of this chapter:
% - Multi-modal dataset with N annotated images across M modalities.
% - SAM-assisted + cross-modal annotation pipeline reducing labeling effort.
% - Baseline detection performance of X mAP@0.5 (RGB) and Y mAP@0.5 (NIR).
% - The trained detector is deployed in Chapter~\ref{chapNBV} as the
%   perception front-end for semantic mapping.
